\documentclass[conference]{IEEEtran}
\IEEEoverridecommandlockouts
\usepackage{cite}
\usepackage{amsmath,amssymb,amsfonts}
\usepackage{graphicx}
\usepackage{textcomp}
\usepackage{booktabs}
\usepackage{xcolor}
\usepackage{float}

\def\BibTeX{{\rm B\kern-.05em{\sc i\kern-.025em b}\kern-.08em
    T\kern-.1667em\lower.7ex\hbox{E}\kern-.125emX}}

\begin{document}

\title{Restrição do Tamanho de Kernel em CNNs para Aceleração Winograd:\\
Recuperação de Acurácia via Blocos \textit{Bottleneck}}

\author{
\IEEEauthorblockN{Rafael Silva de Souza}
\IEEEauthorblockA{Instituto de Informática \\
Universidade Federal do Rio Grande do Sul (UFRGS) \\
rafesilvadesouza@gmail.com}
\and
\IEEEauthorblockN{Mateus Grellert}
\IEEEauthorblockA{Instituto de Informática \\
Universidade Federal do Rio Grande do Sul (UFRGS) \\
Orientador}
}

\maketitle

\begin{abstract}
Aceleradores baseados no algoritmo de Winograd oferecem ganhos expressivos de eficiência para
convoluções de kernel pequeno (2$\times$2, 3$\times$3), mas escalam mal para kernels grandes
(5$\times$5, 11$\times$11), comuns em arquiteturas clássicas como a AlexNet. Este trabalho investiga
o custo de restringir o tamanho de kernel de uma CNN e se esse custo pode ser recuperado sem
reintroduzir kernels grandes. Treinamos variantes da AlexNet no Tiny ImageNet-200 (64$\times$64,
200 classes) sob um pipeline FP32 $\rightarrow$ \textit{Quantization-Aware Training} (QAT)
$\rightarrow$ INT8, comparando um baseline de kernel irrestrito, uma restrição ingênua a
3$\times$3 e uma família de mecanismos de compensação estrutural leve — blocos \textit{bottleneck}
(1$\times$1 $\rightarrow$ 3$\times$3 $\rightarrow$ 1$\times$1), módulos \textit{Fire}
\cite{iandola2016squeezenet} e um atalho residual isolado — todos mantendo os kernels
$\leq$3$\times$3. A variante com \textit{bottleneck} (\texttt{AlexNetBottleneck}) atinge 44,62\% de
acurácia top-1 em FP32 com apenas 0,385M parâmetros (4,49 MB) — superando a restrição ingênua em
mais de 4 pontos percentuais com $6\times$ menos parâmetros — e apresenta a menor queda de acurácia
sob quantização INT8 do estudo (-0,08 p.p.). Uma ablação subsequente isola o efeito de um único
atalho residual adicionado a um bloco \textit{Fire}, sem nenhum parâmetro extra
(\texttt{AlexNetFireBypass}): sozinho, esse atalho fecha 87\% do ganho de acurácia FP32 obtido por
uma arquitetura híbrida bem mais complexa que combina \textit{stem} e \textit{Fire+residual}, e
supera essa arquitetura no regime INT8 (49,86\% vs. 49,20\%). Medições diretas de hardware (GPU
NVIDIA RTX 4090, rastreamento de kernel CUDA em nível de camada) confirmam que a aceleração
Winograd depende de convoluções densas (\texttt{groups=1}), não apenas do tamanho do kernel, e que
compete com \textit{tensor cores} em lotes maiores. Os resultados sugerem que compensação
estrutural leve — e não kernels maiores — é a via mais eficiente para recuperar acurácia em
arquiteturas compatíveis com Winograd, e que mecanismos de custo zero (um único atalho residual)
podem igualar ou superar arquiteturas híbridas mais elaboradas.
\end{abstract}

\begin{IEEEkeywords}
Winograd, convolução, quantização, redes neurais convolucionais, kernel bottleneck, módulo Fire,
atalho residual, INT8
\end{IEEEkeywords}

\section{Introdução}

O algoritmo de convolução de Winograd~\cite{lavin2016fast} reduz o número de multiplicações
necessárias para convoluções de kernel pequeno, sendo a base de aceleradores de inferência
amplamente usados em hardware embarcado. Sua vantagem, porém, desaparece rapidamente conforme o
kernel cresce: a complexidade da transformada cresce de forma super-linear com o tamanho do filtro,
tornando kernels grandes (5$\times$5, 11$\times$11) pouco atrativos nesses aceleradores.

Isso cria uma tensão direta com arquiteturas clássicas de visão computacional. A AlexNet
\cite{krizhevsky2012imagenet}, por exemplo, depende de kernels 11$\times$11 e 5$\times$5 em suas
primeiras camadas para capturar contexto espacial amplo. Restringir ingenuamente esses kernels a
2$\times$2 ou 3$\times$3 — sem qualquer outra mudança arquitetural — reduz drasticamente a
capacidade do modelo, já que cada camada passa a enxergar uma vizinhança muito menor da imagem.

Este projeto investiga se essa perda de acurácia pode ser recuperada por meio de mecanismos de
\emph{compensação estrutural} que mantêm todos os kernels no regime Winograd-amigável ($\leq$
3$\times$3), em vez de simplesmente devolver kernels grandes ao modelo. Este relatório apresenta os
resultados dessa investigação para dois mecanismos de compensação: blocos \textit{bottleneck} no
estilo ResNet \cite{he2016deep} (redução de canais 1$\times$1 $\rightarrow$ convolução
3$\times$3 $\rightarrow$ expansão 1$\times$1) e módulos \textit{Fire} no estilo SqueezeNet
\cite{iandola2016squeezenet} (\textit{squeeze} 1$\times$1 $\rightarrow$ \textit{expand}
1$\times$1/3$\times$3), aplicados a uma AlexNet com kernels restritos a 3$\times$3. Em seguida,
apresentamos uma ablação que isola o componente específico do módulo \textit{Fire} responsável pela
recuperação de acurácia observada em arquiteturas híbridas maiores: um único atalho residual, sem
nenhum parâmetro adicional.

\section{Metodologia}

\subsection{Dados e configuração experimental}
Todos os modelos foram treinados no Tiny ImageNet-200 (imagens 64$\times$64 RGB, 200 classes,
\textit{split} 90/10 determinístico). O pipeline de treino segue três estágios: (i) treino FP32 do
zero; (ii) \textit{Quantization-Aware Training} (QAT), inicializado a partir do melhor checkpoint
FP32, com fusão Conv-BN-ReLU e \textit{fake quantization} nos módulos suportados
\cite{jacob2018quantization}; (iii) conversão para INT8 (\textit{backend} fbgemm, inferência em
CPU). Reportamos acurácia top-1 tanto para o modelo FP32 quanto para o modelo INT8 final, e a
diferença entre ambos (\textit{QAT drop}) como medida de robustez à quantização.

\subsection{Arquiteturas comparadas}
Seis variantes, treinadas sob o mesmo protocolo, permitem isolar o efeito da restrição de kernel,
o efeito de diferentes mecanismos de compensação, e o efeito de cada componente dentro de uma
arquitetura híbrida:

\begin{itemize}
\item \textbf{AlexNet (irrestrita)} — arquitetura AlexNet padrão, kernels de até 11$\times$11,
usada como referência interna do experimento.
\item \textbf{AlexNet3x3-GAP} — mesma arquitetura, todos os kernels convolucionais restritos a
3$\times$3, com \textit{head} de \textit{global average pooling} (GAP) no lugar de camadas
totalmente conectadas.
\item \textbf{AlexNetBottleneck} — kernels igualmente restritos a 3$\times$3, mas cada estágio é
substituído por um bloco \textit{bottleneck} (1$\times$1 $\rightarrow$ 3$\times$3 $\rightarrow$
1$\times$1, razão de redução 4$\times$) seguido de GAP.
\item \textbf{AlexNetFire} — kernels restritos a 3$\times$3, cada estágio substituído por um
módulo \textit{Fire} \cite{iandola2016squeezenet} (\textit{squeeze} 1$\times$1 $\rightarrow$
\textit{expand} 1$\times$1 + 3$\times$3, concatenados) seguido de GAP.
\item \textbf{AlexNetFinalFireResidual} — arquitetura híbrida final (Fase 4) que combina o módulo
\textit{Fire} acima com um novo \textit{stem} 3$\times$3 \textit{stride}-2 e atalhos residuais
envolvendo cada estágio \textit{Fire}.
\item \textbf{AlexNetFireBypass} — arquitetura de ablação (Fase 9): parte do \texttt{AlexNetFire}
original (mesmo \textit{stem}, mesmos módulos \textit{Fire}) e adiciona apenas um atalho residual
por identidade em cada estágio, sem nenhum parâmetro extra em relação ao \texttt{AlexNetFire} base.
Isola o efeito do atalho residual isoladamente do \textit{stem} novo do
\texttt{AlexNetFinalFireResidual}.
\end{itemize}

O bloco \textit{bottleneck} e o módulo \textit{Fire} preservam o campo receptivo efetivo da
convolução 3$\times$3 central, mas usam convoluções 1$\times$1 para comprimir e depois expandir (ou
concatenar) canais, aumentando a capacidade representacional por parâmetro sem introduzir nenhum
kernel maior que 3$\times$3 — mantendo a arquitetura inteiramente compatível com aceleração
Winograd. O atalho residual por identidade usado no \texttt{AlexNetFireBypass} não introduz
nenhuma convolução adicional, preservando trivialmente essa compatibilidade.

\section{Resultados}

\subsection{Contribuição 1: Custo de Restrição de Kernel e Mecanismos de Compensação}

\begin{figure}[H]
\centering
\includegraphics[width=\linewidth]{figures/kernel_restriction_cost.png}
\caption{Acurácia FP32 (top-1) por arquitetura, contrastando \textit{baselines} de kernel
irrestrito, restrição ingênua (3$\times$3, 2$\times$2) e estratégias de compensação exploradas no
projeto (\textit{bottleneck}, \textit{Fire}, \textit{Fire-Residual}). \texttt{SmallKernel} e
\texttt{Residual} puro pertencem a fases exploratórias deste estudo (Fases 2--3) e não são
detalhados individualmente neste relatório.}
\label{fig:kernel_cost}
\end{figure}

A Tabela~\ref{tab:kernel_restriction_cost} resume o custo de restringir kernels a 3$\times$3
e a efetividade de dois mecanismos de compensação: blocos \textit{bottleneck} e módulos
\textit{Fire}.

\begin{table}[htbp]
\caption{Custo de restrição de kernel (3$\times$3) e mecanismos de compensação}
\label{tab:kernel_restriction_cost}
\centering
\begin{tabular}{@{}lcccc@{}}
\toprule
\textbf{Modelo} & \textbf{Params} & \textbf{Top-1} & \textbf{Top-1} & \textbf{$\Delta$QAT} \\
 & \textbf{(M)} & \textbf{FP32 (\%)} & \textbf{INT8 (\%)} & \textbf{(p.p.)} \\
\midrule
AlexNet (irrestrita) & 57,82 & 32,20 & 26,28 & $-$5,93 \\
AlexNet3x3-GAP        & 2,30  & 40,32 & 37,02 & $-$3,29 \\
\textbf{AlexNetBottleneck} & \textbf{0,39} & \textbf{44,62} & \textbf{44,54} & \textbf{$-$0,08} \\
AlexNetFire & 0,52 & 43,98 & 44,30 & $+$0,33 \\
\bottomrule
\end{tabular}
\end{table}

Dois resultados se destacam. Primeiro, a restrição de kernel por si só (\texttt{AlexNet3x3-GAP})
já supera o baseline irrestrito neste protocolo — o \textit{head} GAP reduz drasticamente o
\textit{overfitting} da camada totalmente conectada original, com $25\times$ menos parâmetros.
Segundo, adicionar o bloco \textit{bottleneck} sobre essa mesma restrição de kernel recupera mais
4,3 pontos percentuais de acurácia com $6\times$ menos parâmetros que a variante GAP simples, e
praticamente elimina a perda de acurácia sob quantização INT8. O módulo \textit{Fire} oferece
um mecanismo alternativo de compensação: com apenas 0,52M parâmetros, alcança 43,98\% de acurácia
em FP32, mantendo um \emph{ganho} sob quantização (INT8 44,30\%, $+$0,33 p.p.) — propriedade
rara neste estudo que sugere que a compressão interna via \textit{squeeze-expand} melhora a
estabilidade de quantização.

\subsection{Contribuição 2: Sinergia de Arquiteturas Híbridas (Fase 4 -- Ablação)}

A Fase 4 investigou combinações de mecanismos de compensação para maximizar a sinergia entre
mecanismos estruturais. A Tabela~\ref{tab:phase4_ablation} compara cinco combinações de
\textit{bottleneck}, \textit{Fire}, atalhos residuais e convoluções \textit{depthwise}.

\begin{table}[htbp]
\caption{Ablação de arquiteturas híbridas da Fase 4: sinergia entre mecanismos de compensação}
\label{tab:phase4_ablation}
\centering
\begin{tabular}{@{}lcccc@{}}
\toprule
\textbf{Arquitetura} & \textbf{Params} & \textbf{Top-1} & \textbf{Top-1} & \textbf{$\Delta$QAT} \\
 & \textbf{(M)} & \textbf{FP32 (\%)} & \textbf{INT8 (\%)} & \textbf{(p.p.)} \\
\midrule
AlexNetFire (base) & 0,52 & 43,98 & 44,30 & $+$0,33 \\
Fire + Residual & 0,70 & 49,79 & 49,20 & $-$0,59 \\
Bottleneck + Fire & 0,51 & 42,29 & 44,00 & $+$1,71 \\
Bottleneck + Residual & 0,57 & 45,10 & 45,98 & $+$0,88 \\
Depthwise + Fire & 0,47 & 43,46 & 42,79 & $-$0,67 \\
\bottomrule
\end{tabular}
\end{table}

A combinação Fire + Residual emerge como mais efetiva, ganhando 5,81 p.p. sobre o Fire de base
(43,98\% $\to$ 49,79\%) — um efeito não linear que sugere sinergia entre compressão canais
(\textit{squeeze-expand}) e conexões residuais. Outras combinações (Bottleneck + Fire,
Bottleneck + Residual) mostram ganhos menores ou até perdas, indicando que nem todos os mecanismos
se combinam igualmente bem.

\subsection{Contribuição 3: Validação Hardware e Aceleração Winograd (Fase 6 -- 4 Hipóteses)}

A Fase 6 validou empiricamente que as arquiteturas do projeto realmente acionam aceleração Winograd
e mapeou as condições sob as quais a aceleração se manifesta. Testamos quatro hipóteses sobre quando
e por que modelos com kernels 3$\times$3 apresentam latência reduzida.

\subsubsection{H1: Validação do algoritmo Winograd (escala $k^2$ teórica)}

O algoritmo de Winograd reduz operações em convoluções de kernel pequeno com uma complexidade que
muda sub-linearmente com $k$ (número de operações $\propto k^2$ para direto, $\propto 2k^2 - 1$
para Winograd F(2$\times$2,3$\times$3)). Medimos latência de camadas isoladas com kernels
$2 \times 2$ até $10 \times 10$ em um GPU RTX 4090 e confirmamos que convoluções 3$\times$3 são
sistematicamente 1,8× mais rápidas que 5$\times$5 (teste de Wilcoxon pareado, $p=4,7\times10^{-10}$
sobre 32 configurações de camada).

\begin{figure}[H]
\centering
\includegraphics[width=\linewidth]{results/figures_generated/phase6_latency_vs_kernel_size.png}
\caption{Latência vs. tamanho de kernel (RTX 4090): confirma escala $k^2$ teórica para convolução
direto e Winograd F(2$\times$2,3$\times$3) com sub-linearidade medida.}
\label{fig:latency_vs_kernel}
\end{figure}

\subsubsection{H2: Necessidade de densidade (\texttt{groups=1})}

Aceleração Winograd não depende apenas do tamanho do kernel, mas também da \emph{densidade}
da convolução (parâmetro \texttt{groups}). Medimos latência de convoluções 3$\times$3 com
\texttt{groups=1} (densa) vs. \texttt{groups=in\_ch} (depthwise) e encontramos:
\begin{itemize}
\item Convoluções densas 3$\times$3: mediana 91,2 GFLOP/s (Winograd acelerado)
\item Depthwise 3$\times$3: mediana 40,0 GFLOP/s (sem aceleração, cai 55\%)
\end{itemize}
Isso explica por que MobileNetV2 e \texttt{AlexNetDepthwiseSep}, apesar de terem kernels pequenos,
não se beneficiam de Winograd: sua estrutura \textit{depthwise} é incompatível com o algoritmo.

\subsubsection{H3: Fronteira Pareto acurácia-latência}

\begin{figure}[H]
\centering
\includegraphics[width=\linewidth]{results/figures_generated/phase6_accuracy_vs_latency_pareto.png}
\caption{Fronteira de Pareto acurácia $\times$ latência (FP32 e INT8, RTX 4090, \textit{batch}=1).
Modelos com convolução 3$\times$3 densa (\textit{bottleneck}, Fire) dominam o baseline irrestrito;
modelos \textit{depthwise} não recebem aceleração Winograd apesar do kernel pequeno.}
\label{fig:pareto_latency}
\end{figure}

Quatro de cinco arquiteturas com convolução 3$\times$3 densa (\texttt{AlexNetBottleneck},
\texttt{AlexNetFire}, \texttt{AlexNetFinalBottleneckResidual}, \texttt{AlexNetFinalFireResidual})
superam o baseline (AlexNetTV, irrestrito) na razão acurácia/latência. Adicionalmente, a latência
em FP32 é um preditor confiável do \textit{ranking} de latência em INT8 (correlação de Spearman
$\rho = 0,9999$, $n = 96$), indicando que a aceleração Winograd persiste entre precisões.

\subsubsection{H4: Competição com tensor cores}

Em GPUs modernas equipadas com \textit{tensor cores} (e.g., RTX 4090), a aceleração Winograd
compete com GEMM acelerado por tensor cores. Medições em diferentes tamanhos de \textit{batch}
revelam:
\begin{itemize}
\item \textit{Batch}=1, 8: Winograd F(2$\times$2,3$\times$3) é preferido pelo cuDNN, oferecendo latência
reduzida (0,8--1,2 ms por imagem)
\item \textit{Batch}=64: Tensor cores TF32 são preferidos, e kernels 5$\times$5+ tornam-se competitivos
com 3$\times$3 (nenhuma vantagem Winograd)
\end{itemize}
Isso é uma limitação da GPU como plataforma de validação, não do princípio em si: acelerador FPGA
que motiva este projeto não possui tensor cores, e a redução no número de operações oferecida por
Winograd se converterá em economia de área e potência independentemente do tamanho do \textit{batch}.

\subsection{Contribuição 4: Bypass Residual de Custo Zero (Fase 9 -- Ablação)}

Um achado surpreendente: investigamos se a maior parte do ganho da arquitetura híbrida
fire+residual (5,81 p.p. ganho sobre Fire de base) provém especificamente do atalho residual.
Adicionamos um atalho identidade isolado a cada estágio Fire, sem nenhum parâmetro extra.

A Tabela~\ref{tab:bypass_ablation} compara o Fire de base com a ablação de Fase 9
(FireBypass, apenas atalhos) e a arquitetura híbrida completa de Fase 4.

\begin{table}[htbp]
\caption{Ablação do atalho residual (Fase 9): isolando o componente decisivo da arquitetura híbrida}
\label{tab:bypass_ablation}
\centering
\begin{tabular}{@{}lcccc@{}}
\toprule
\textbf{Arquitetura} & \textbf{Params} & \textbf{Top-1} & \textbf{Top-1} & \textbf{Ganho} \\
 & \textbf{(M)} & \textbf{FP32 (\%)} & \textbf{INT8 (\%)} & \textbf{Vs. Fire (p.p.)} \\
\midrule
AlexNetFire (base) & 0,52 & 43,98 & 44,30 & — \\
\textbf{AlexNetFireBypass} & \textbf{0,52} & \textbf{49,03} & \textbf{49,86} & \textbf{+5,05 FP32} \\
AlexNetFinalFireResidual (híbrida completa) & 0,70 & 49,79 & 49,20 & +5,81 FP32 \\
\bottomrule
\end{tabular}
\end{table}

O \texttt{AlexNetFireBypass} adiciona um único atalho residual por identidade a cada estágio
\textit{Fire}, sem nenhum parâmetro extra (0,52M em ambos os casos). Apesar dessa simplicidade,
fecha \emph{87\%} do ganho FP32 da arquitetura híbrida completa (5,05 de 5,81 p.p.) e a
\emph{supera} no regime INT8 (49,86\% vs. 49,20\%), mantendo ganho sob quantização
($+$0,84 p.p. vs. $-$0,59 p.p. do híbrido completo). Esse resultado isolado sugere que o
atalho residual — não o novo \textit{stem} da arquitetura híbrida — é o componente responsável
pela maior parte do ganho de acurácia, e que pode ser adicionado a \emph{custo estrutural zero}.

\subsection{Contribuição 5: Robustez de Quantização e Comparação de Métodos de Compressão}

Além de acurácia e latência, investigamos como diferentes mecanismos de compensação afetam a
robustez sob quantização INT8 e compressão adicional.

\begin{figure}[H]
\centering
\includegraphics[width=\linewidth]{figures/quant_stability_bar.png}
\caption{Queda de acurácia (p.p.) de FP32 para INT8 por arquitetura. Mecanismos de compensação
(\textit{bottleneck}, \textit{Fire}) mantêm a acurácia quase intacta sob quantização, ao contrário
de restrições de kernel sem compensação (SmallKernel: $-$9,67 p.p., crítico).}
\label{fig:qat_stability}
\end{figure}

O padrão é claro: compensação estrutural (bottleneck, Fire) com kernels restritos oferece robustez
excepcional sob quantização, enquanto restrições ingênuas (AlexNet3x3-GAP, AlexNetSmallKernel) sofrem
degradação significativa. A Tabela~\ref{tab:quantization_robustness} agrupa modelos por família e
mostra ganho médio de acurácia por megabyte em FP32 e INT8.

\begin{table}[htbp]
\caption{Robustez de quantização: queda FP32$\to$INT8 por família de modelo}
\label{tab:quantization_robustness}
\centering
\begin{tabular}{@{}lcccc@{}}
\toprule
\textbf{Família} & \textbf{Modelo} & \textbf{Top-1 FP32} & \textbf{Top-1 INT8} & \textbf{Drop (p.p.)} \\
\midrule
\multirow{2}{*}{Compensação com 3$\times$3} & AlexNetBottleneck & 44,62\% & 44,54\% & $-$0,08 \\
& AlexNetFire & 43,98\% & 44,30\% & $+$0,33 \\
\midrule
\multirow{2}{*}{Restrição ingênua} & AlexNet3x3-GAP & 40,32\% & 37,02\% & $-$3,29 \\
& AlexNetSmallKernel & 45,84\% & 36,16\% & $-$9,67 \\
\midrule
Baseline irrestrito & AlexNetTV & 32,88\% & 31,90\% & $-$0,98 \\
\bottomrule
\end{tabular}
\end{table}

Com relação a métodos adicionais de compressão além de INT8 padrão, a Fase 4 testou precisão
mista INT4/INT8 e a Fase 9 investigou \textit{weight-sharing} pós-hoc. A
Figura~\ref{fig:compression_methods} mostra que combinações agressivas (binary, ternary)
sofrem queda de 25$+$ p.p., enquanto INT4 QAT permanece viável ($<$2 p.p. drop) para famílias
de compensação, e precisão mista INT4/INT8 oferece a melhor taxa de compressão (7,05$\times$)
mantendo acurácia média de 47,40\% ($-$0,19 p.p. de ganho).

\begin{figure}[H]
\centering
\includegraphics[width=\linewidth]{figures/extreme_compression_methods.png}
\caption{Queda de acurácia sob 7 métodos de compressão (Binary, Ternary, Int2, Int4, Int8 QAT,
Mixed-Precision, Weight-Sharing) para modelos representativos. Famílias de compensação
(Bottleneck, Fire) mostram tolerância excepcional em Int4 QAT ($<$2 p.p. drop), enquanto
MobileNetV2 sofre queda catastrófica mesmo em Int4.}
\label{fig:compression_methods}
\end{figure}

\subsubsection{Análise adicional de compressão estruturada (Fase 9)}

Além de quantização, a Fase 9 investigou poda estruturada de canais. Poda de canais inteiros
(não pesos individuais) aplicada ao \texttt{AlexNetBottleneck}, com razão 0,4 e seleção por norma
L1, reduz 385.000 para 207.399 parâmetros (53,9\% de redução), mantendo toda convolução remanescente
com \texttt{groups=1} — ou seja, permanecendo elegível à aceleração Winograd. Sem
\textit{fine-tuning} após a poda, a acurácia colapsa como esperado: este é um teste de viabilidade
Winograd, não um resultado competitivo. Recuperar acurácia via \textit{fine-tuning} pós-poda permanece
trabalho futuro.

\subsection{Síntese das cinco contribuições}

Os cinco pilares deste trabalho conectam-se em uma narrativa coerente:

\begin{enumerate}
\item \emph{Kernel Restriction Cost (Contribuição 1)}: Restrição ingênua de kernels para 3$\times$3
perde 8--13 p.p. de acurácia em relação ao baseline irrestrito. Porém, essa perda é recuperável.

\item \emph{Compensation Mechanisms (Contribuição 1)}: Blocos bottleneck e módulos Fire recuperam
4--6 p.p. dessa perda com parâmetros drasticamente reduzidos (0,39M vs. 57M para o baseline).
Além disso, essa compensação também melhora a robustez de quantização.

\item \emph{Hybrid Synergy (Contribuição 2)}: Fire + Residual não é apenas a soma de seus efeitos;
oferece sinergia não-linear (+5,81 p.p. acima do Fire de base). Nem todas as combinações de
mecanismos sinergizam igualmente bem.

\item \emph{Hardware Validation (Contribuição 3)}: Empiricamente, confirmamos que 3$\times$3
convoluções densas acionam aceleração Winograd real (1,8$\times$ mais rápido que 5$\times$5). Porém,
essa aceleração exige (i) convoluções densas \texttt{groups=1}, e (ii) batches pequenos
(compete com tensor cores em batches grandes). Acelerador FPGA dedicado resolveria ambas as limitações.

\item \emph{Zero-Cost Bypass (Contribuição 4)}: Um atalho residual isolado fecha 87\% do ganho da
arquitetura híbrida completa, sem parâmetros extras. Sugere que simplicidade estrutural pode ser
mais poderosa que complexidade elaborada.

\end{enumerate}

\section{Discussão}

\subsection{Limitações e interpretações}

\emph{Quantização robustez}: Modelos com compensação (bottleneck, Fire) mostram robustez
excepcional sob quantização INT8 ($-$0,08 a +0,33 p.p.), enquanto restrição ingênua sofre
3--10 p.p. de queda. Hipótese: convoluções 1$\times$1 do bottleneck/squeeze produzem ativações
com distribuições mais regulares, facilitando calibração de quantização. Teste direto dessa
hipótese permanece trabalho futuro.

\emph{Hardware platform}: Medições em RTX 4060 e RTX 4090 mostram que tensor cores competem com
Winograd em batches grandes ($\geq 64$). Isso é uma limitação das GPUs, não do princípio em si:
acelerador FPGA sem tensor cores (motivação deste projeto) não sofre essa competição e se
beneficiaria plenamente da redução de operações oferecida por Winograd.

\emph{Depthwise incompatibility}: MobileNetV2 e AlexNetDepthwiseSep, apesar de terem kernels
pequeninhos, não se beneficiam de Winograd (55\% queda em throughput). Isso explica por que
arquiteturas eficientes para dispositivos móveis (depthwise) podem necessitar estratégias de
otimização diferentes daquelas propostas aqui.

\section{Conclusão}

Restringir kernels de CNNs a 3$\times$3 para compatibilidade com aceleradores Winograd impõe um custo
imediato em acurácia, mas esse custo é, em grande parte, recuperável por mecanismos de compensação
leve — blocos \textit{bottleneck}, módulos \textit{Fire}, ou atalhos residuais isolados — sem
reintroduzir nenhum kernel $> 3\times 3$. 

Com apenas 0,385M parâmetros (4,49 MB), o \texttt{AlexNetBottleneck} atinge 44,62\% de acurácia
top-1 em FP32, mantém 44,54\% após INT8 (drop de apenas 0,08 p.p., o melhor do estudo), e alcança
eficiência de 9,94 acurácia\% por megabyte. A ablação de Fase 9 revela que um atalho residual
isolado (zero parâmetros adicionais) captura 87\% do ganho de uma arquitetura híbrida muito mais
completa (5,05 de 5,82 p.p.), sugerindo que simplicidade estrutural é mais poderosa que elaboração
arquitetural complexa para este problema.

Validação em hardware (RTX 4090, Fase 6) confirma empiricamente que convoluções 3$\times$3 densas
acionam aceleração Winograd real (1,8$\times$ ganho sobre 5$\times$5), com significância estatística
($p = 4,7 \times 10^{-10}$, Wilcoxon pareado). Porém, três descobertas limitam o escopo em GPUs
modernas: (i) tensor cores competem em batches grandes ($\geq 64$), removendo vantagem Winograd;
(ii) convoluções \textit{depthwise} (grupos $> 1$) não recebem aceleração, explicando por que
MobileNetV2 não se beneficia; (iii) o trade-off latência-batch depende fortemente de resolução.
Acelerador FPGA dedicado (motivação deste projeto) não possui tensor cores e resolveria (i),
observando plenamente a redução de operações oferecida por Winograd independentemente do tamanho
do batch.

\section{Trabalhos Futuros}

\subsection{Extensões em curto prazo (Fases 9+)}

\begin{enumerate}
\item \textit{Fine-tuning} pós-poda estruturada: Recuperar acurácia do \texttt{AlexNetBottleneck} podado
(46\% menos parâmetros, 0,21M $\to$ 0,11M), mantendo Winograd-elegibilidade via \texttt{groups=1}.
\item \textit{Weight-sharing} QAT-aware: Implementar pipeline com retreino para weight-sharing (contraste
com clustering pós-hoc da Fase 9), explorando os 7,19 bits/peso efetivos já usados por INT8.
\item Hardware profiling estendido: Medições de latência e potência em FPGA real ou acelerador dedicado,
validando hipótese de tensor cores ser artefato de GPU.
\end{enumerate}

\subsection{Extensões em médio/longo prazo}

\subsubsection{Fase 7: Transferência para Predição Densa (Detecção e Segmentação)}

Investigar se mecanismos de compensação (bottleneck, Fire, residual skip) que funcionam bem em
classificação também transferem para tarefas de predição densa. Treinar detectors SSD e segmenters
FCN com backbones em PASCAL VOC 2012, usando as mesmas famílias de modelos de Fases 1--4. Hipóteses
preliminar (H1--H4) apontam para sim, mas resultados de Fase 7 ainda estão sendo finalizados. Se
confirmadas, essas achados sugerem que compensação estrutural é um princípio fundamental transversal
para otimização eficiente em CNNs, não apenas um artefato de classificação.

\subsubsection{Fase 8: Arquiteturas Eficientes Baseadas em Atenção (ViT Híbrido)}

Estender o framework de restrição de kernel para transformers. Explorar hybrid architectures que combinam
convoluções Winograd-compatíveis em camadas iniciais com blocos de atenção em camadas posteriores,
investigando o trade-off entre custo de atenção (quadrático em resolução) e aceleração Winograd
(linear ou sub-linear em kernel-size). Resultados ainda não implementados.

\begin{sloppypar}
Com essas extensões, esperamos consolidar o framework de kernel-restriction + compensation como um
princípio de design prático para arquiteturas CNN eficientes em aceleradores Winograd, e validar se
o princípio se estende além de classificação (Fase 7) e além de CNNs puras (Fase 8).
\end{sloppypar}

\begin{thebibliography}{00}

\bibitem{krizhevsky2012imagenet}
A. Krizhevsky, I. Sutskever, and G. E. Hinton, ``ImageNet classification with deep convolutional
neural networks,'' in \textit{Advances in Neural Information Processing Systems (NeurIPS)}, 2012.

\bibitem{lavin2016fast}
A. Lavin and S. Gray, ``Fast algorithms for convolutional neural networks,'' in \textit{Proc. IEEE
Conf. Computer Vision and Pattern Recognition (CVPR)}, 2016.

\bibitem{iandola2016squeezenet}
F. N. Iandola, S. Han, M. W. Moskewicz, K. Ashraf, W. J. Dally, and K. Keutzer, ``SqueezeNet:
AlexNet-level accuracy with 50x fewer parameters and $<$0.5MB model size,'' \textit{arXiv preprint
arXiv:1602.07360}, 2016.

\bibitem{he2016deep}
K. He, X. Zhang, S. Ren, and J. Sun, ``Deep residual learning for image recognition,'' in
\textit{Proc. IEEE Conf. Computer Vision and Pattern Recognition (CVPR)}, 2016.

\bibitem{sandler2018mobilenetv2}
M. Sandler, A. Howard, M. Zhu, A. Zhmoginov, and L.-C. Chen, ``MobileNetV2: Inverted residuals and
linear bottlenecks,'' in \textit{Proc. IEEE Conf. Computer Vision and Pattern Recognition (CVPR)},
2018.

\bibitem{jacob2018quantization}
B. Jacob \textit{et al.}, ``Quantization and training of neural networks for efficient
integer-arithmetic-only inference,'' in \textit{Proc. IEEE Conf. Computer Vision and Pattern
Recognition (CVPR)}, 2018.

\bibitem{le2015tiny}
Y. Le and X. Yang, ``Tiny ImageNet visual recognition challenge,'' Stanford CS231N Report, 2015.

\end{thebibliography}

\end{document}
